% GUID v6, Part III: the one place where the scoring rules of Parts II and III are fixed (pull, threshold, classes, BIC, bootstrap,
% multiple tests, preregistration). The tally itself is Section sec:tally.
\section{How the confrontations are scored}\label{sec:scoring}

Sections~\ref{sec:data} to~\ref{sec:exotics} confront Part~I with measurements by one set of rules, fixed here; Primer~\ref{box:stats} gives
the statistics behind them.

% GUID v6, Primer P8 (label box:stats): the statistics behind the scoring rules of Section sec:scoring. Background only; it states
% no result of the paper and does not repeat the rules themselves. Place it in Section sec:scoring: % GUID v6, Primer P8 (label box:stats): the statistics behind the scoring rules of Section sec:scoring. Background only; it states
% no result of the paper and does not repeat the rules themselves. Place it in Section sec:scoring: % GUID v6, Primer P8 (label box:stats): the statistics behind the scoring rules of Section sec:scoring. Background only; it states
% no result of the paper and does not repeat the rules themselves. Place it in Section sec:scoring: \input{primers/p8-stats}
\begin{primer}{Statistics}{box:stats}
\emph{Pulls and 95\%.} A measurement $x\pm\sigma$ usually means a Gaussian error with standard deviation $\sigma$. If a prediction
$x_G\pm\sigma_G$ is right and the two errors are independent, the difference $x-x_G$ is Gaussian with variance $\sigma^2+\sigma_G^2$, so the
\emph{pull} $p=(x-x_G)/\sqrt{\sigma^2+\sigma_G^2}$ is a standard normal number. It falls outside $\pm1.96$ with probability 5\% (two-sided),
so $|p|<1.96$ is a test at 95\%, and a published 95\% half-width $b$ corresponds to $\sigma=b/1.96$. \emph{Worked pull:} a model predicts
$10.0\pm0.3$ and an experiment finds $11.0\pm0.4$. Then $p=1.0/\sqrt{0.4^2+0.3^2}=1.0/0.5=2.0$: a FAIL at 95\%, although a right model
gives $|p|\ge2.0$ in 4.6\% of independent tests. For $n$ independent tests, about $0.05\,n$ FAILs are expected by chance, and the sum of the
squared pulls follows a $\chi^2$ distribution with $n$ degrees of freedom.

\emph{BIC and $\Delta$BIC.} A model with $k$ fitted parameters that reaches the maximum likelihood $\mathcal L_{\max}$ on $N$ data points has
$\mathrm{BIC}=-2\ln\mathcal L_{\max}+k\ln N$~\cite{schwarz1978}. The penalty $k\ln N$ charges each extra parameter. For large $N$ and
independent points, the difference $\Delta$BIC of two models approximates $-2$ times the logarithm of their Bayes factor, the ratio of the
probabilities of the data under the two models. A $\Delta$BIC of 2 to 6 is then positive evidence, 6 to 10 strong and above 10 very
strong~\cite{kassraftery1995}. The approximation assumes $N$ independent points. If the points come in groups that share an error, such
as the distance of one galaxy, the effective number of independent data is closer to the number of groups. A point-wise $\Delta$BIC then
overstates the evidence, roughly in proportion to the number of points per group.

\emph{The bootstrap.} To learn how much a fitted value would vary in a new sample, one draws many new samples of the same size from the
data themselves, with replacement, and refits each~\cite{efron1979}. The spread of the refitted values estimates the uncertainty; the 16th
to 84th percentiles bracket a range comparable to $\pm1\sigma$. The units resampled must be the independent ones, for rotation curves whole
galaxies, not single points.

\emph{Many tests.} With $m$ tests at 5\% each, the chance of at least one false positive, the \emph{family-wise error rate}, grows towards
$1-0.95^m$. Bonferroni's remedy tests each at $0.05/m$. Holm's method orders the $p$-values, $p_{(1)}\le\dots\le p_{(m)}$, and rejects
$p_{(i)}$ while $p_{(i)}\le0.05/(m-i+1)$; it keeps the family-wise error rate at 5\% under any dependence and never rejects less than
Bonferroni~\cite{holm1979}. Fisher's method combines $n$ \emph{independent} $p$-values: under the null each $p_i$ is uniform, so
$-2\ln p_i$ follows $\chi^2$ with 2 degrees of freedom, and $X^2=-2\sum_i\ln p_i$ follows $\chi^2$ with $2n$~\cite{fisher1932}. It responds
strongly to a few very small $p_i$.

\emph{Look-elsewhere and preregistration.} A search over many masses, scales or candidate values finds a small \emph{local} $p$-value
somewhere by chance. The \emph{global} $p$-value, the chance of so small a value anywhere in the search, is larger, roughly the local one
times a trials factor~\cite{grossvitells2010}. A test is confirmatory only if its hypothesis, data, statistic and decision rule were
fixed before the outcome was seen; \emph{preregistration} writes them down in a dated record beforehand, and every choice made afterwards
is then visible as exploratory~\cite{nosek2018}.
\end{primer}

\begin{primer}{Statistics}{box:stats}
\emph{Pulls and 95\%.} A measurement $x\pm\sigma$ usually means a Gaussian error with standard deviation $\sigma$. If a prediction
$x_G\pm\sigma_G$ is right and the two errors are independent, the difference $x-x_G$ is Gaussian with variance $\sigma^2+\sigma_G^2$, so the
\emph{pull} $p=(x-x_G)/\sqrt{\sigma^2+\sigma_G^2}$ is a standard normal number. It falls outside $\pm1.96$ with probability 5\% (two-sided),
so $|p|<1.96$ is a test at 95\%, and a published 95\% half-width $b$ corresponds to $\sigma=b/1.96$. \emph{Worked pull:} a model predicts
$10.0\pm0.3$ and an experiment finds $11.0\pm0.4$. Then $p=1.0/\sqrt{0.4^2+0.3^2}=1.0/0.5=2.0$: a FAIL at 95\%, although a right model
gives $|p|\ge2.0$ in 4.6\% of independent tests. For $n$ independent tests, about $0.05\,n$ FAILs are expected by chance, and the sum of the
squared pulls follows a $\chi^2$ distribution with $n$ degrees of freedom.

\emph{BIC and $\Delta$BIC.} A model with $k$ fitted parameters that reaches the maximum likelihood $\mathcal L_{\max}$ on $N$ data points has
$\mathrm{BIC}=-2\ln\mathcal L_{\max}+k\ln N$~\cite{schwarz1978}. The penalty $k\ln N$ charges each extra parameter. For large $N$ and
independent points, the difference $\Delta$BIC of two models approximates $-2$ times the logarithm of their Bayes factor, the ratio of the
probabilities of the data under the two models. A $\Delta$BIC of 2 to 6 is then positive evidence, 6 to 10 strong and above 10 very
strong~\cite{kassraftery1995}. The approximation assumes $N$ independent points. If the points come in groups that share an error, such
as the distance of one galaxy, the effective number of independent data is closer to the number of groups. A point-wise $\Delta$BIC then
overstates the evidence, roughly in proportion to the number of points per group.

\emph{The bootstrap.} To learn how much a fitted value would vary in a new sample, one draws many new samples of the same size from the
data themselves, with replacement, and refits each~\cite{efron1979}. The spread of the refitted values estimates the uncertainty; the 16th
to 84th percentiles bracket a range comparable to $\pm1\sigma$. The units resampled must be the independent ones, for rotation curves whole
galaxies, not single points.

\emph{Many tests.} With $m$ tests at 5\% each, the chance of at least one false positive, the \emph{family-wise error rate}, grows towards
$1-0.95^m$. Bonferroni's remedy tests each at $0.05/m$. Holm's method orders the $p$-values, $p_{(1)}\le\dots\le p_{(m)}$, and rejects
$p_{(i)}$ while $p_{(i)}\le0.05/(m-i+1)$; it keeps the family-wise error rate at 5\% under any dependence and never rejects less than
Bonferroni~\cite{holm1979}. Fisher's method combines $n$ \emph{independent} $p$-values: under the null each $p_i$ is uniform, so
$-2\ln p_i$ follows $\chi^2$ with 2 degrees of freedom, and $X^2=-2\sum_i\ln p_i$ follows $\chi^2$ with $2n$~\cite{fisher1932}. It responds
strongly to a few very small $p_i$.

\emph{Look-elsewhere and preregistration.} A search over many masses, scales or candidate values finds a small \emph{local} $p$-value
somewhere by chance. The \emph{global} $p$-value, the chance of so small a value anywhere in the search, is larger, roughly the local one
times a trials factor~\cite{grossvitells2010}. A test is confirmatory only if its hypothesis, data, statistic and decision rule were
fixed before the outcome was seen; \emph{preregistration} writes them down in a dated record beforehand, and every choice made afterwards
is then visible as exploratory~\cite{nosek2018}.
\end{primer}

\begin{primer}{Statistics}{box:stats}
\emph{Pulls and 95\%.} A measurement $x\pm\sigma$ usually means a Gaussian error with standard deviation $\sigma$. If a prediction
$x_G\pm\sigma_G$ is right and the two errors are independent, the difference $x-x_G$ is Gaussian with variance $\sigma^2+\sigma_G^2$, so the
\emph{pull} $p=(x-x_G)/\sqrt{\sigma^2+\sigma_G^2}$ is a standard normal number. It falls outside $\pm1.96$ with probability 5\% (two-sided),
so $|p|<1.96$ is a test at 95\%, and a published 95\% half-width $b$ corresponds to $\sigma=b/1.96$. \emph{Worked pull:} a model predicts
$10.0\pm0.3$ and an experiment finds $11.0\pm0.4$. Then $p=1.0/\sqrt{0.4^2+0.3^2}=1.0/0.5=2.0$: a FAIL at 95\%, although a right model
gives $|p|\ge2.0$ in 4.6\% of independent tests. For $n$ independent tests, about $0.05\,n$ FAILs are expected by chance, and the sum of the
squared pulls follows a $\chi^2$ distribution with $n$ degrees of freedom.

\emph{BIC and $\Delta$BIC.} A model with $k$ fitted parameters that reaches the maximum likelihood $\mathcal L_{\max}$ on $N$ data points has
$\mathrm{BIC}=-2\ln\mathcal L_{\max}+k\ln N$~\cite{schwarz1978}. The penalty $k\ln N$ charges each extra parameter. For large $N$ and
independent points, the difference $\Delta$BIC of two models approximates $-2$ times the logarithm of their Bayes factor, the ratio of the
probabilities of the data under the two models. A $\Delta$BIC of 2 to 6 is then positive evidence, 6 to 10 strong and above 10 very
strong~\cite{kassraftery1995}. The approximation assumes $N$ independent points. If the points come in groups that share an error, such
as the distance of one galaxy, the effective number of independent data is closer to the number of groups. A point-wise $\Delta$BIC then
overstates the evidence, roughly in proportion to the number of points per group.

\emph{The bootstrap.} To learn how much a fitted value would vary in a new sample, one draws many new samples of the same size from the
data themselves, with replacement, and refits each~\cite{efron1979}. The spread of the refitted values estimates the uncertainty; the 16th
to 84th percentiles bracket a range comparable to $\pm1\sigma$. The units resampled must be the independent ones, for rotation curves whole
galaxies, not single points.

\emph{Many tests.} With $m$ tests at 5\% each, the chance of at least one false positive, the \emph{family-wise error rate}, grows towards
$1-0.95^m$. Bonferroni's remedy tests each at $0.05/m$. Holm's method orders the $p$-values, $p_{(1)}\le\dots\le p_{(m)}$, and rejects
$p_{(i)}$ while $p_{(i)}\le0.05/(m-i+1)$; it keeps the family-wise error rate at 5\% under any dependence and never rejects less than
Bonferroni~\cite{holm1979}. Fisher's method combines $n$ \emph{independent} $p$-values: under the null each $p_i$ is uniform, so
$-2\ln p_i$ follows $\chi^2$ with 2 degrees of freedom, and $X^2=-2\sum_i\ln p_i$ follows $\chi^2$ with $2n$~\cite{fisher1932}. It responds
strongly to a few very small $p_i$.

\emph{Look-elsewhere and preregistration.} A search over many masses, scales or candidate values finds a small \emph{local} $p$-value
somewhere by chance. The \emph{global} $p$-value, the chance of so small a value anywhere in the search, is larger, roughly the local one
times a trials factor~\cite{grossvitells2010}. A test is confirmatory only if its hypothesis, data, statistic and decision rule were
fixed before the outcome was seen; \emph{preregistration} writes them down in a dated record beforehand, and every choice made afterwards
is then visible as exploratory~\cite{nosek2018}.
\end{primer}


\paragraph{The pull and the 95\% rule.} A confrontation sets a GUID value $x_G$ against a datum. Where it exists, a theory uncertainty
$\sigma_G$ comes with $x_G$. It is propagated from measured inputs, such as Planck's $\Omega_m$, or it is a model's own error, such as the
scatter of a fitted mass formula. For a Gaussian datum $x\pm\sigma$ (one standard deviation) the \emph{pull} is
\begin{equation}\label{eq:pull}
p=\frac{x-x_G}{\sqrt{\sigma^2+\sigma_G^2}} .
\end{equation}
A published 95\% half-width $b$ enters as $\sigma=b/1.96$; for an asymmetric error, the side towards the prediction is used. A confrontation
is \textbf{OK} when $|p|<1.96$, that is at 95\% for a Gaussian, and a \textbf{FAIL} otherwise. The same threshold holds in every section.
Other data are allowed regions, probabilities or qualitative statements of the experimenters. Such a confrontation is OK when $x_G$ lies
inside the region. Where a region is quoted at 90\% CL the row says so; passing a 90\% region implies passing at 95\%. A FAIL stays a FAIL
when general relativity, $\Lambda$CDM or the Standard Model share it. Its verdict word is then \vd{inherited} (Section~\ref{sec:howto}).

\paragraph{The classes.} Every row has one class. A row that is not scored states the reason.
\begin{itemize}[nosep,leftmargin=*]
\item \emph{Scored}: a prediction confronted with a datum it was not set from, with the verdict OK or FAIL. Only scored rows enter the tally
(Section~\ref{sec:tally}).
\item \emph{Not yet sensitive}: a scored row whose prediction is specific to GUID, while the datum lies far from where GUID departs from
general relativity, $\Lambda$CDM, the Standard Model or quantum mechanics. ``Far'' means that the departure is smaller than the datum's
reach by a factor of ten or more, or that it scales with a free parameter that may be zero. Such a row can still fail, so it stays scored.
When it holds, it shows only that GUID's departure is small.
\item \emph{Conditional}: the same test for a model that Part~I uses but GUID does not derive, such as a plateau inflation potential.
It is reported apart and not counted for GUID.
\item \emph{Consistency} (CONS): the prediction was set from this datum, so it holds by construction. It is never counted as a success.
\item \emph{Information} (INFO): superseded by a newer measurement that includes its data, derived from scored rows, overlapping a scored
row, inconclusive, a check of a theory GUID adopts unchanged, or without a usable uncertainty. It is reported, not scored.
\end{itemize}
A scored row is also either \emph{shared}, when its prediction is that of general relativity, $\Lambda$CDM, the Standard Model or quantum
mechanics, or \emph{GUID-specific}. Only a GUID-specific row that is sensitive can tell GUID apart from the theories it adopts.

\paragraph{Chance and shared data.} If GUID were right and $n$ scored data were independent, about $0.05n$ FAILs would occur by chance.
Data often share inputs, such as the microwave background, so FAILs can cluster. Each section names the rows that share data.

\paragraph{Likelihoods and BIC.} Fits of whole curves, such as rotation curves, compare models by their likelihood $\mathcal L$. For a model
with $k$ fitted parameters and $N$ data points, the Bayesian information criterion~\cite{schwarz1978} is
\begin{equation}\label{eq:bic}
\mathrm{BIC}=-2\ln\mathcal L_{\max}+k\ln N .
\end{equation}
The model with the lower BIC is preferred. $\Delta$BIC is the difference to a reference model, negative when the model fits better. For
independent points, a difference above 10 counts as very strong evidence~\cite{kassraftery1995}. The points of one galaxy are not
independent, however: they share its distance, inclination and stellar mass-to-light ratio. A point-wise $\Delta$BIC counts each galaxy's
common error once per point and so overstates the evidence. A parameter fitted per galaxy also absorbs these shared errors, which favours
any model that has one. Sections~\ref{sec:sparc} and~\ref{sec:remnant} therefore report $\Delta$BIC but base their verdicts on the galaxy
bootstrap and on galaxy-by-galaxy counts.

\paragraph{The galaxy bootstrap.} Ranges of fitted parameters come from a bootstrap over galaxies~\cite{efron1979}. The sample is redrawn
with replacement, whole galaxies at a time, and refitted; 1000 draws use a fixed random seed. A range in brackets is the 16th to 84th
percentile. Resampling galaxies rather than points keeps each galaxy's correlated errors together.

\paragraph{Many tests.} When many hypotheses are tested at once, some pass the 95\% threshold by chance.
\begin{itemize}[nosep,leftmargin=*]
\item The \emph{Holm correction}~\cite{holm1979} keeps the chance of any false positive in a family of $m$ tests below 5\%. It orders the
$p$-values, compares the smallest with $0.05/m$, the next with $0.05/(m-1)$, and so on, and stops at the first that is not small enough.
\item \emph{Fisher's method}~\cite{fisher1932} combines $n$ independent $p$-values into $X^2=-2\sum_i\ln p_i$. Under the null hypothesis
$X^2$ follows a $\chi^2$ distribution with $2n$ degrees of freedom.
\item The \emph{look-elsewhere effect}~\cite{grossvitells2010} is the chance of finding a coincidence somewhere when many places or
candidate values are tried. We account for it by computing the probability of a match at least as close over every candidate tried
(Sections~\ref{sec:remnant} and~\ref{sec:exotics}).
\end{itemize}

\paragraph{Preregistration.} A prediction fixed before the data are analysed cannot be tuned to them. The Standard Model statements and
their scoring were committed in \texttt{docs/\allowbreak simulation/\allowbreak sm\_preregistration.md} before any analysis (Section~\ref{sec:sm}). The candidate
scales and history proxies of Section~\ref{sec:remnant}, and the assignment of model levels to exotic hadrons in Section~\ref{sec:exotics},
were fixed before the comparison. Each section lists, as ``choices made after seeing the data'', what it decided later and whether a verdict
depends on it.

\paragraph{Sources.} Every measured number was checked against the arXiv text or abstract of its source. Rows read only through secondary
sources are listed in each section's caveats.
